\chapterbackmatter{Solutions to Supplementary Exercises}

\begin{enumerate}
  \SupplementarySolution{1}{The Zinn--Justin Identity and Linearized Cohomology}
  \begin{enumerate}
    \item[(a)]
    With a constant Grassmann parameter \(\theta\), make the change of
    variables
    \(\chi^n\mapsto\chi^n+\theta\Delta^n\).  The action is invariant,
    and the \(K\)-term changes by
    \(K_n s\Delta^n=K_ns^2\chi^n=0\).  Invariance of the measure and the
    absence of a surface term in field space therefore give
    \[
      0=\int d^4x\,J_n(x)
        \left.\frac{\delta_R W}{\delta K_n(x)}\right|_J .
    \]
    At fixed \(K\), the Legendre transform gives
    \[
      \frac{\delta_L\Gamma}{\delta\chi^n(x)}=-J_n(x),
      \qquad
      \frac{\delta_R\Gamma}{\delta K_n(x)}
        =\left.\frac{\delta_R W}{\delta K_n(x)}\right|_J .
    \]
    Substitution yields
    \[
      \boxed{\int d^4x\,
        \frac{\delta_R\Gamma}{\delta K_n}
        \frac{\delta_L\Gamma}{\delta\chi^n}=0.}
    \]
    The assumptions are BRST invariance of the gauge-fixed action,
    nilpotence, an invariant regularized measure, and boundary conditions
    that permit a change of integration variables.  A nontrivial Jacobian
    is the possible anomaly.
    \item[(b)]
    Since \(s\) carries dimension one, ghost number one, and odd parity,
    \[
      [s\chi^n]=d_n+1,\quad
      {\rm gh}(s\chi^n)=\gamma_n+1,\quad
      \epsilon(s\chi^n)=\epsilon_n+1\pmod2 .
    \]
    Requiring \(\int K_ns\chi^n\) to be even, dimensionless, and of ghost
    number zero gives
    \[
      \boxed{[K_n]=3-d_n,\quad
      {\rm gh}(K_n)=-\gamma_n-1,\quad
      \epsilon(K_n)=\epsilon_n+1\pmod2 .}
    \]
    The resulting table is
    \[
    \begin{array}{c|ccc}
      \chi^n&[K_n]&{\rm gh}(K_n)&\epsilon(K_n)\\ \hline
      A_\mu&2&-1&1\\
      \omega&2&-2&0\\
      \omega^*&2&0&0\\
      \psi&3/2&-1&0\\
      \phi&2&-1&1
    \end{array}
    \]
    where \(\psi\) is fermionic and \(\phi\) is bosonic.  Although the
    formal rule would assign a source of dimension one and ghost number
    \(-1\) to \(h\), no such source is useful because \(sh=0\); it would
    couple identically to zero.
    \item[(c)]
    The graded Jacobi identity gives, for even \(S\),
    \[
      (S,(S,X))=\frac12((S,S),X).
    \]
    Hence \((S,S)=0\) implies
    \(\mathcal B_S^2X=0\).  Now deform
    \(S\mapsto S+\varepsilon\Sigma\), with \(\varepsilon^2=0\).  Then
    \[
      (S+\varepsilon\Sigma,S+\varepsilon\Sigma)
        =(S,S)+2\varepsilon(S,\Sigma)
        =2\varepsilon\mathcal B_S\Sigma .
    \]
    The master equation is preserved to first order exactly when
    \(\mathcal B_S\Sigma=0\).  Thus admissible infinitesimal counterterms
    are cocycles of the linearized Slavnov operator.
  \end{enumerate}

  \SupplementarySolution{2}{Loopwise Master Equations and Trivial Counterterms}
  \begin{enumerate}
    \item[(a)]
    Substitution of \(\Gamma=\sum_N\Gamma_N\) into
    \((\Gamma,\Gamma)=0\) gives at total loop order \(N\)
    \[
      \boxed{\sum_{r=0}^{N}(\Gamma_r,\Gamma_{N-r})=0.}
    \]
    Suppose all lower-order divergences and subdivergences have been
    canceled.  In the pole part of this equation, every term with
    \(0<r<N\) is finite.  Only the two equal endpoint terms can contain
    the remaining overall divergence:
    \[
      (S_R,\Gamma_{N,\infty})
        +(\Gamma_{N,\infty},S_R)=0.
    \]
    Since both entries are even, the two terms are equal in the chapter's
    antibracket convention.  Therefore
    \[
      \boxed{(S_R,\Gamma_{N,\infty})=0.}
    \]
    Locality of \(\Gamma_{N,\infty}\) follows from the subtraction of all
    subdivergences.
    \item[(b)]
    An odd functional \(F\) generates the infinitesimal anticanonical
    transformation
    \[
      \delta\chi^n=(\chi^n,F),\qquad
      \delta K_n=(K_n,F).
    \]
    The induced change of \(S\) is
    \[
      \delta S=(S,F)=\mathcal B_SF=\Sigma .
    \]
    Anticanonical transformations preserve the antibracket, so the master
    equation remains true.  If \(F\) is linear in the antifields it is an
    ordinary local field redefinition, with the compensating inverse
    transformation of sources.  If \(F\) is a change of gauge-fixing
    fermion, its effect is \(s\,\delta\Psi\).  In either case the cocycle is
    exact: it changes coordinates or gauge choice, not the coefficient of a
    nontrivial gauge-invariant operator.
    \item[(c)]
    A constant shift gives
    \[
      \boxed{\int d^4x\,
        \frac{\delta\Gamma}{\delta\omega_\alpha^*(x)}=0,}
    \]
    so every local dependence on \(\omega_\alpha^*\) must be through a
    derivative.  Independently, the source equation for the linear
    transformation \(s\omega_\alpha^*=-h_\alpha\) is exact:
    \[
      \frac{\delta_L\Gamma}{\delta K_\alpha^*(x)}
        =-h_\alpha(x).
    \]
    It has no radiative correction because the inserted operator is linear.
    Power counting and ghost number exclude all quadratic source terms
    except, provisionally, \(K_\alpha^*K_\beta^*\): each \(K^*\) has
    dimension two and ghost number zero.  But such a term would make
    \(\delta\Gamma/\delta K_\alpha^*\) depend on \(K^*\), contradicting the
    exact source equation.  Therefore the divergent effective action is at
    most linear in every BRST source.
  \end{enumerate}

  \SupplementarySolution{3}{Scalar BRST Cohomology and Power Counting}
  \begin{enumerate}
    \item[(a)]
    The undifferentiated \(A^2\) term makes the ghost action depend on
    \(\omega^*\) without a derivative, so a constant antighost shift is no
    longer a symmetry.  Since
    \[
      [\omega^*A^2]=3,\quad
      {\rm gh}(\omega^*A^2)=-1,
    \]
    its BRST variation is an allowed dimension-four, ghost-number-zero
    interaction.  Likewise
    \[
      [\omega^*\omega^*\omega]=3,\qquad
      {\rm gh}(\omega^*\omega^*\omega)=-1,
    \]
    so
    \[
    \begin{split}
      s(\omega_\alpha^*\omega_\beta^*\omega_\gamma)
        ={}&-h_\alpha\omega_\beta^*\omega_\gamma
          +\omega_\alpha^*h_\beta\omega_\gamma\\
         &-\frac12C_{\gamma\delta\epsilon}
          \omega_\alpha^*\omega_\beta^*\omega_\delta\omega_\epsilon
    \end{split}
    \]
    is also allowed.  The last term is a four-ghost vertex absent from the
    minimal Faddeev--Popov action, yet no remaining symmetry forbids its
    divergent coefficient.  Closure under renormalization therefore
    requires starting from the general BRST-invariant action
    \(I_0+s\Psi\), with this interaction available as a counterterm.
    \item[(b)]
    Using the odd Leibniz rule,
    \[
    \begin{split}
      s^2\phi
        &=i(s\omega^\alpha)T_\alpha\phi
          -i\omega^\alpha T_\alpha(s\phi)\\
        &=-\frac{i}{2}C_{\alpha\beta\gamma}
            \omega^\beta\omega^\gamma T_\alpha\phi
          +\omega^\alpha\omega^\beta T_\alpha T_\beta\phi\\
        &=\frac12\omega^\alpha\omega^\beta
          \left([T_\alpha,T_\beta]
            -iC_{\gamma\alpha\beta}T_\gamma\right)\phi .
    \end{split}
    \]
    Thus \(s^2\phi=0\) for arbitrary ghosts precisely when
    \([T_\alpha,T_\beta]=iC_{\gamma\alpha\beta}T_\gamma\).  Conversely, if
    nilpotence holds on a faithful collection of scalar fields, the
    coefficient of each independent
    \(\omega^\alpha\omega^\beta\) vanishes and gives the representation
    relation.
    \item[(c)]
    The scalar source has dimension two, ghost number \(-1\), and odd
    parity.  Two such sources already have dimension four, so the only
    candidate is
    \(K_{\phi i}K_{\phi j}\), with an antisymmetric coefficient if the
    sources are distinct.  Its ghost number is \(-2\), so it is forbidden.
    The mixed dimension-four candidates with other dimension-two sources
    are
    \[
      K_\phi K_A,\qquad K_\phi K_\omega,\qquad
      K_\phi K_{\omega^*}.
    \]
    Their ghost numbers are respectively \(-2,-3,-1\).  Adding a ghost to
    repair the last value would raise the dimension to five.  The product
    \(K_\phi K_\psi\) has dimension \(7/2\), ghost number \(-2\), and an
    uncontracted spinor index; a field that contracts it makes the dimension
    exceed four.  Hence no local Lorentz scalar of dimension at most four
    and ghost number zero is quadratic in \(K_\phi\) or mixes \(K_\phi\)
    with another BRST source.
  \end{enumerate}

  \SupplementarySolution{4}{Slavnov Stability under Renormalization}
  Write the broken identity at the first nonvanishing order as
  \[
    \frac12(\Gamma,\Gamma)=\hbar^N\mathcal A_N+O(\hbar^{N+1}).
  \]
  Applying the antibracket with \(\Gamma\) and using the Jacobi identity
  gives \((\Gamma,(\Gamma,\Gamma))=0\).  At order \(N\), lower orders may
  be replaced by \(S_R\), so
  \[
    \boxed{\mathcal B_{S_R}\mathcal A_N
      =(S_R,\mathcal A_N)=0.}
  \]
  If \(\mathcal A_N=\mathcal B_{S_R}X_N\) for a local ghost-number-zero
  functional \(X_N\), adding the counterterm \(-\hbar^NX_N\) restores the
  identity.  A closed but non-exact ghost-number-one local functional is a
  gauge anomaly.  Its cohomology class, rather than a regulator-dependent
  representative, is the obstruction.

  \SupplementarySolution{5}{Invariant Scalar and Yukawa Counterterms}
  \begin{enumerate}
    \item[(a)]
    Write
    \[
      V(\phi)=a_i\phi_i+\frac12m^2_{ij}\phi_i\phi_j
        +\frac1{3!}\kappa_{ijk}\phi_i\phi_j\phi_k
        +\frac1{4!}\lambda_{ijkl}\phi_i\phi_j\phi_k\phi_l .
    \]
    With \(\delta\phi_i=i\epsilon^\alpha
    (T_\alpha)_i{}^j\phi_j\), invariance gives
    \[
    \begin{aligned}
      &a_i(T_\alpha)_i{}^j=0,\\
      &(T_\alpha)_i{}^{i'}m^2_{i'j}
        +(T_\alpha)_j{}^{j'}m^2_{ij'}=0,\\
      &(T_\alpha)_i{}^{i'}\kappa_{i'jk}
        +(T_\alpha)_j{}^{j'}\kappa_{ij'k}
        +(T_\alpha)_k{}^{k'}\kappa_{ijk'}=0,
    \end{aligned}
    \]
    and the analogous four-term identity for \(\lambda\).  Under
    \(\phi'=M\phi\), the generators transform by similarity and each
    coupling by the corresponding tensor power of \(M^{-1}\).  Every
    displayed zero therefore maps to zero.  The invariant tensors, modulo
    these basis changes, are exactly the possible scalar-potential
    counterterms.
    \item[(b)]
    The variations are
    \[
      \delta\psi_R=i\epsilon^\alpha t_\alpha^R\psi_R,\quad
      \delta\bar\psi_L=-i\epsilon^\alpha\bar\psi_Lt_\alpha^L,\quad
      \delta\phi_i=i\epsilon^\alpha(T_\alpha)_i{}^j\phi_j .
    \]
    The coefficient of each \(\epsilon^\alpha\phi_j\) vanishes when
    \[
      \boxed{
      t_\alpha^LY^j-Y^jt_\alpha^R
        -(T_\alpha)_i{}^jY^i=0.}
    \]
    This says that \(Y\) is an invariant map from the scalar representation
    tensored with the right-handed representation to the left-handed one.
    A divergent local Yukawa vertex has dimension four and the same
    external fields.  The ghost-number-zero Slavnov identity reduces on
    that vertex to the same infinitesimal gauge-invariance equation, so its
    tensor lies in the same invariant subspace and can be absorbed into
    \(Y^i\).
  \end{enumerate}

  \SupplementarySolution{6}{Background Splitting, Gauge Fixing, and Determinants}
  \begin{enumerate}
    \item[(a)]
    In matrix notation, with \(D_\mu[A]\epsilon
    =\partial_\mu\epsilon+[A_\mu,\epsilon]\), take
    \[
    \begin{array}{c|cc}
       &\delta A_\mu&\delta A'_\mu\\ \hline
      \text{background}
        &D_\mu[A]\epsilon&[A'_\mu,\epsilon]\\
      \text{quantum}
        &0&D_\mu[A+A']\epsilon .
    \end{array}
    \]
    The background transformation gives
    \[
      \delta(A_\mu+A'_\mu)
        =\partial_\mu\epsilon+[A_\mu+A'_\mu,\epsilon]
        =D_\mu[A+A']\epsilon ,
    \]
    which is exactly the variation of the sum under the quantum
    transformation.  The two operations distribute the same total
    transformation differently between the prescribed and integrated
    fields.
    \item[(b)]
    Since an adjoint field satisfies
    \(\delta A'_\mu=[A'_\mu,\epsilon]\), covariance of the background
    derivative gives
    \[
      \delta(\bar D^\mu A'_\mu)
        =[\bar D^\mu A'_\mu,\epsilon].
    \]
    In chapter components this is
    \[
      \delta f_\alpha
        =-C_{\alpha\beta\gamma}\epsilon_\beta f_\gamma .
    \]
    Therefore
    \[
      \delta(f_\alpha f_\alpha)
        =-2C_{\alpha\beta\gamma}
          f_\alpha\epsilon_\beta f_\gamma=0,
    \]
    because the two \(f\)'s are symmetric under
    \(\alpha\leftrightarrow\gamma\), whereas the structure constants are
    antisymmetric.
    \item[(c)]
    The quantum transformation is
    \[
      \delta A'_{\alpha\mu}
        =\bar D_\mu\epsilon_\alpha
         -C_{\alpha\beta\gamma}\epsilon_\beta A'_{\gamma\mu}.
    \]
    Hence
    \[
      \delta f_\alpha
        =\bar D^\mu\!\left[
          \bar D_\mu\epsilon_\alpha
          -C_{\alpha\beta\gamma}\epsilon_\beta A'_{\gamma\mu}\right].
    \]
    Replacing \(\epsilon\) by \(\omega\) gives
    \[
      \mathcal L_{\rm gh}
        =\omega_\alpha^*\bar D^\mu\!\left[
          \bar D_\mu\omega_\alpha
          -C_{\alpha\beta\gamma}\omega_\beta A'_{\gamma\mu}\right],
    \]
    or, after a background-covariant integration by parts,
    \[
      \mathcal L_{\rm gh}
        =-(\bar D^\mu\omega_\alpha^*)
         \left[\bar D_\mu\omega_\alpha
         -C_{\alpha\beta\gamma}\omega_\beta A'_{\gamma\mu}\right].
    \]
    These are Eqs.~\eqref{eq:17.4.22} and \eqref{eq:17.4.23} with
    background ghost fields set to zero.
    \item[(d)]
    Use
    \[
      sA'_{\alpha\mu}
        =\bar D_\mu\omega_\alpha
         -C_{\alpha\beta\gamma}\omega_\beta A'_{\gamma\mu},\quad
      s\omega_\alpha^*=-h_\alpha,\quad sh_\alpha=0 .
    \]
    The odd Leibniz rule yields
    \[
    \begin{split}
      s\Psi=\int d^4x\Bigl\{&
        h_\alpha\bar D^\mu A'_{\alpha\mu}
        +\frac{\xi}{2}h_\alpha h_\alpha\\
        &+\omega_\alpha^*\bar D^\mu[
          \bar D_\mu\omega_\alpha
          -C_{\alpha\beta\gamma}\omega_\beta A'_{\gamma\mu}]
        \Bigr\}.
    \end{split}
    \]
    The algebraic equation
    \(h_\alpha=-\bar D^\mu A'_{\alpha\mu}/\xi\) produces
    \(-f_\alpha f_\alpha/(2\xi)\), plus the ghost term found above.
    Nilpotence is off shell because \(h\) was retained before elimination.
    Every object in braces contracts adjoint indices covariantly, so the
    expression is background gauge invariant.
    \item[(e)]
    For an infinitesimal parameter \(\epsilon_\alpha(x)\), invariance of the
    action and measure gives
    \[
      \boxed{\bar D_\mu^{\alpha\beta}[A]\,
         \frac{\delta\Gamma[A]}{\delta A_\mu^\beta}=0.}
    \]
    Differentiate once with respect to \(A_\nu^\gamma(y)\) and then set
    \(A=0\).  The term obtained by differentiating \(\bar D\) is
    proportional to the one-point function, which vanishes in a symmetric
    vacuum.  Thus
    \[
      \partial_x^\mu\Gamma_{\mu\nu}^{\alpha\gamma}(x-y)=0,
      \qquad
      q^\mu\Gamma_{\mu\nu}^{\alpha\gamma}(q)=0 .
    \]
    Lorentz and global gauge invariance consequently give
    \[
      \Gamma_{\mu\nu}^{\alpha\beta}(q)
        =\delta^{\alpha\beta}
         (q^2\eta_{\mu\nu}-q_\mu q_\nu)\Pi(q^2).
    \]
    \item[(f)]
    For a real commuting vector \(x\), a complex Grassmann pair
    \(\bar\eta,\eta\), and a quadratic matrix \(M\),
    \[
      \int d^nx\,e^{-\frac12x^{\mathsf T}Mx}
        \propto(\det M)^{-1/2},\qquad
      \int d^n\bar\eta\,d^n\eta\,e^{-\bar\eta M\eta}
        =\det M .
    \]
    Taking logarithms, a real boson contributes
    \(+\tfrac12\operatorname{Tr}\ln M\) to the Euclidean one-loop effective
    action, whereas a complex Grassmann field contributes
    \(-\operatorname{Tr}\ln M\).  A Dirac fermion and a Faddeev--Popov ghost
    pair are both complex Grassmann systems, but their operators and
    multiplicities differ.  Thus the gauge fluctuation supplies the
    \(+\tfrac12\) trace, while fermion and ghost loops carry the closed
    Grassmann-loop minus sign.
  \end{enumerate}

  \SupplementarySolution{7}{Adjoint Scalars, Gauge Covariance, and BRST Gauge Independence}
  \begin{enumerate}
    \item[(a)]
    Write
    \[
      A'_\mu=UA_\mu U^\dagger+\frac{i}{g}U\partial_\mu U^\dagger,
      \qquad \phi'=U\phi U^\dagger.
    \]
    Expanding \(D'_\mu\phi'\) and using
    \(\partial_\mu U^\dagger=-U^\dagger(\partial_\mu U)U^\dagger\), the
    two terms containing \(\partial_\mu U\) cancel, leaving
    \[
      \boxed{D'_\mu\phi'=U(D_\mu\phi)U^\dagger.}
    \]
    Acting twice and antisymmetrizing gives
    \[
    \begin{split}
      [D_\mu,D_\nu]\phi
       &=-ig[\partial_\mu A_\nu-\partial_\nu A_\mu,\phi]
         -g^2\bigl([A_\mu,[A_\nu,\phi]]
                  -[A_\nu,[A_\mu,\phi]]\bigr)\\
       &=-ig[F_{\mu\nu},\phi],
    \end{split}
    \]
    where the Jacobi identity was used in the last step.  Covariance of the
    left-hand side implies
    \[
      [F'_{\mu\nu},U\phi U^\dagger]
        =[UF_{\mu\nu}U^\dagger,U\phi U^\dagger]
    \]
    for arbitrary adjoint \(\phi\), hence
    \(F'_{\mu\nu}=UF_{\mu\nu}U^\dagger\).  Cyclicity of the trace now makes
    \(\operatorname{Tr}F^2\), \(\operatorname{Tr}(D\phi)^2\), and
    \(\operatorname{Tr}\phi^3\) separately invariant.
    \item[(b)]
    The commutator and anticommutator give
    \[
      T_aT_b=\frac1{2N}\delta_{ab}\mathbf1
        +\frac12(d_{ab}{}^c+if_{ab}{}^c)T_c,
    \]
    and therefore
    \[
      \operatorname{Tr}(T_aT_bT_c)
        =\frac14(d_{abc}+if_{abc}).
    \]
    Since \(\phi^a\phi^b\phi^c\) is totally symmetric, its contraction with
    the antisymmetric \(f_{abc}\) vanishes.  Thus the entire interaction is
    \[
      \boxed{-\frac1{3!}\operatorname{Tr}(\phi^3)
        =-\frac1{24}d_{abc}\phi^a\phi^b\phi^c.}
    \]
    This also shows explicitly that it vanishes for \(SU(2)\), for which
    \(d_{abc}=0\).
    \item[(c)]
    \begin{enumerate}
      \item[(i)] Put \(U=1+ig\vartheta\omega\), where \(\vartheta\) is a
      constant Grassmann-odd parameter.  The coefficient of \(\vartheta\)
      in the finite transformations is
      \[
        \boxed{sA_\mu=D_\mu\omega
          =\partial_\mu\omega-ig[A_\mu,\omega],\qquad
        s\phi=ig[\omega,\phi].}
      \]
      Together with \(s\omega=ig\omega^2\), these transformations square to
      zero.  For example, the graded Leibniz rule gives
      \[
      \begin{split}
        s^2\phi
          &=ig\bigl([s\omega,\phi]-\{\omega,s\phi\}\bigr)\\
          &=-g^2[\omega^2,\phi]
            +g^2\{\omega,[\omega,\phi]\}=0,
      \end{split}
      \]
      because \(\{\omega,[\omega,\phi]\}=[\omega^2,\phi]\).  The
      calculation for \(A_\mu\) is the covariant-derivative version of the
      same identity.
      \item[(ii)] Gauge invariance gives \(s\int\mathcal L=0\).  Nilpotence
      gives \(s(s\Psi)=0\), so
      \[
        \boxed{sS_\Psi=0.}
      \]
      This statement is off shell because the auxiliary field satisfies
      \(s\omega^*=-h\) and \(sh=0\).
      \item[(iii)] For a product \({\cal O}\) of gauge-invariant operators,
      \(s{\cal O}=0\).  Differentiating its normalized expectation value
      under \(\Psi\to\Psi+\lambda\delta\Psi\) yields
      \[
        \frac{d}{d\lambda}\langle{\cal O}\rangle_\lambda\Big|_{0}
        =i\bigl(\langle{\cal O}\,s\delta\Psi\rangle
          -\langle{\cal O}\rangle\langle s\delta\Psi\rangle\bigr).
      \]
      BRST invariance of the measure implies \(\langle sX\rangle=0\).
      Taking \(X=\delta\Psi\) kills the second term, while
      \(s({\cal O}\delta\Psi)=(-1)^{|{\cal O}|}{\cal O}s\delta\Psi\)
      kills the first.  Hence the derivative vanishes.  A BRST anomaly or
      a nonzero boundary term in field space is precisely what would spoil
      the conclusion.
    \end{enumerate}
  \end{enumerate}

  \SupplementarySolution{8}{Fundamental and Adjoint Matter in a Gluon Self-Energy}
  \begin{enumerate}
    \item[(a)]
    Direct substitution of
    \[
      A'_\mu=UA_\mu U^\dagger+\frac{i}{g}U\partial_\mu U^\dagger
    \]
    into the definition of \(F_{\mu\nu}\) cancels all terms with two
    derivatives of \(U\), as well as the cross terms with one derivative.
    The result is
    \[
      \boxed{F'_{\mu\nu}=UF_{\mu\nu}U^\dagger.}
    \]
    For \(U=1+i\alpha^cT_c+O(\alpha^2)\),
    \[
      \delta F_{\mu\nu}
       =i[\alpha^cT_c,F_{\mu\nu}^bT_b]
       =f_{bc}{}^aF_{\mu\nu}^b\alpha^cT_a,
    \]
    where
    \[
      \boxed{[T_b,T_c]=if_{bc}{}^aT_a.}
    \]
    Hence \(\delta F_{\mu\nu}^a=f_{bc}{}^aF_{\mu\nu}^b\alpha^c\), exactly
    as stated.  Raising or lowering adjoint indices with \(\delta_{ab}\)
    makes \(f_{abc}\) totally antisymmetric.
    \item[(b)]
    With \(\operatorname{Tr}(T_aT_b)=\delta_{ab}/2\), the source coupling is
    \(J_a^\mu A_\mu^a=2\operatorname{Tr}(J^\mu A_\mu)\).  Its homogeneous
    part is invariant only when
    \[
      \boxed{J^\mu\longmapsto UJ^\mu U^\dagger.}
    \]
    More concretely, define
    \[
      D_\mu\psi=(\partial_\mu-igA_\mu^aT_a)\psi .
    \]
    The inhomogeneous term in \(A'_\mu\) gives
    \(D'_\mu\psi'=U D_\mu\psi\), so the invariant fermion Lagrangian is
    \(-\overline\psi(\gamma^\mu D_\mu+m)\psi\).  Expanding it and comparing
    with \(J_a^\mu A_\mu^a\) gives
    \[
      \boxed{J_a^\mu=ig\,\overline\psi\gamma^\mu T_a\psi.}
    \]
    The factor \(i\) is a consequence of the volume's mostly-plus gamma
    convention; converting back to the source's mostly-minus matrices
    yields the familiar \(g\overline\psi\Gamma^\mu T_a\psi\).
    The bilinear transforms in the adjoint because
    \(\overline\psi\mapsto\overline\psi U^\dagger\).
    \item[(c)]
    Infinitesimally,
    \[
      \delta\Psi=i[\alpha^aT_a,\Psi^bT_b]
        =f_{bc}{}^a\Psi^b\alpha^cT_a,
      \qquad
      \boxed{\delta\Psi^a=f_{bc}{}^a\Psi^b\alpha^c.}
    \]
    The adjoint covariant derivative and its transformation are
    \[
      D_\mu\Psi=\partial_\mu\Psi-ig[A_\mu,\Psi],
      \qquad D'_\mu\Psi'=U(D_\mu\Psi)U^\dagger.
    \]
    In the trace normalization used in the question, one gauge-invariant
    Lagrangian is therefore
    \[
      \boxed{\mathcal L_{\rm adj}
       =-\frac14\operatorname{Tr}(F_{\mu\nu}F^{\mu\nu})
        -\operatorname{Tr}\!\left[
          \overline\Psi(\gamma^\mu D_\mu+m)\Psi\right].}
    \]
    Overall trace normalizations can of course be absorbed into the
    corresponding coefficients without changing the gauge-invariance
    argument.
    \item[(d)]
    The required graph list is
    \begin{center}
    % The tadpole loop is an explicit closed Bezier.  Drawn as
    % "to[out=110,in=70,looseness=5] (v)" it starts and ends at the same
    % point, so the |to| path computes a zero-length chord, divides by it and
    % aborts the picture with "Dimension too large".  Loop arrows are placed
    % mid-arc so that no arrowhead lands on a gluon attachment point, and each
    % scope carries vertex blobs.
    \begin{tikzpicture}[scale=0.8,
      gluon/.style={decorate,decoration={snake,amplitude=1.1pt,segment length=4pt},thick},
      flow/.style={thick,postaction={decorate},decoration={markings,
        mark=at position 0.5 with {\arrow{Stealth[length=4.4pt,width=3.2pt]}}}},
      every node/.style={font=\scriptsize}]
      \begin{scope}[shift={(0,1.25)}]
        \draw[gluon] (-1,0)--(1,0); \node at (0,-.9) {tree};
      \end{scope}
      \begin{scope}[shift={(3,1.25)}]
        \coordinate (l) at (-.55,0); \coordinate (r) at (.55,0);
        \draw[gluon] (-1.25,0)--(l); \draw[gluon] (r)--(1.25,0);
        \draw[flow] (l) arc (180:0:.55);
        \draw[flow] (r) arc (0:-180:.55);
        \fill (l) circle (1.4pt); \fill (r) circle (1.4pt);
        \node at (0,-.9) {fermion bubble};
      \end{scope}
      \begin{scope}[shift={(6,1.25)}]
        \coordinate (l) at (-.55,0); \coordinate (r) at (.55,0);
        \draw[gluon] (-1.25,0)--(l); \draw[gluon] (r)--(1.25,0);
        \draw[flow,dashed] (l) arc (180:0:.55);
        \draw[flow,dashed] (r) arc (0:-180:.55);
        \fill (l) circle (1.4pt); \fill (r) circle (1.4pt);
        \node at (0,-.9) {ghost bubble};
      \end{scope}
      \begin{scope}[shift={(1.5,-1.25)}]
        \coordinate (l) at (-.55,0); \coordinate (r) at (.55,0);
        \draw[gluon] (-1.25,0)--(l); \draw[gluon] (r)--(1.25,0);
        \draw[gluon] (l) arc (180:0:.55);
        \draw[gluon] (r) arc (0:-180:.55);
        \fill (l) circle (1.4pt); \fill (r) circle (1.4pt);
        \node at (0,-.9) {gluon bubble};
      \end{scope}
      \begin{scope}[shift={(4.8,-1.25)}]
        \coordinate (v) at (0,0);
        \draw[gluon] (-1.25,0)--(v); \draw[gluon] (v)--(1.25,0);
        \draw[gluon] (v) .. controls (-.58,.32) and (-.58,.98) .. (0,1.06)
                          .. controls (.58,.98) and (.58,.32) .. (v);
        \fill (v) circle (1.4pt); \node at (0,-.9) {four-gluon tadpole};
      \end{scope}
    \end{tikzpicture}
    \end{center}
    The fermion graph has color trace
    \[
      \operatorname{tr}_r(T_aT_b)=T(r)\delta_{ab}.
    \]
    All pure-gauge and ghost graphs are unchanged when the matter
    representation changes.  For \(SU(N)\),
    \[
      T(\boldsymbol N)=\frac12,\qquad
      T(\mathrm{adj})=C_A=N.
    \]
    Thus one adjoint Dirac fermion multiplies the fermion-loop contribution
    by
    \[
      \boxed{\frac{T(\mathrm{adj})}{T(\boldsymbol N)}=2N}
    \]
    relative to one fundamental Dirac fermion.  An adjoint Majorana fermion
    has half the Dirac multiplicity.
  \end{enumerate}
  \SupplementarySolution{9}{A Complex Scalar in the Gauge Beta Function}
  Expanding the covariant kinetic term gives
  \[
  \begin{split}
    (D_\mu\Phi)^\dagger D^\mu\Phi
      ={}&\partial_\mu\Phi^\dagger\partial^\mu\Phi
      +igA_\mu^a\Phi^\dagger T^a
        \overleftrightarrow{\partial^\mu}\Phi\\
      &+g^2A_\mu^aA^{b\mu}\Phi^\dagger T^aT^b\Phi .
  \end{split}
  \]
  Thus the cubic vertex is proportional to
  \(gT^a(k+k')_\mu\), while the seagull is proportional to
  \(g^2\eta_{\mu\nu}\{T^a,T^b\}\).  With a common overall \(i\) fixed by
  the Feynman-rule convention, the two diagrams combine as
  \[
    i\Pi_{\mu\nu}^{ab}(p)=
    g^2T(r)\delta^{ab}\mu^{2\epsilon}
    \int\frac{d^dk}{(2\pi)^d}
    \left[
      \frac{(2k+p)_\mu(2k+p)_\nu}{k^2(k+p)^2}
      -\frac{2\eta_{\mu\nu}}{k^2}
    \right].
  \]
  The second integral vanishes when the scalar is massless because it has
  no scale.  It is nevertheless required by the gauge-invariant
  Lagrangian and cancels the corresponding local term if a mass or another
  regulator is retained.

  For the bubble, introduce a Feynman parameter and shift
  \(k=\ell-xp\).  Symmetric integration replaces
  \(\ell_\mu\ell_\nu\) by
  \(\eta_{\mu\nu}\ell^2/d\).  The pole part is transverse:
  \[
    \Pi_{\mu\nu}^{ab}(p)\big|_{\rm pole}
      =\delta^{ab}(p^2\eta_{\mu\nu}-p_\mu p_\nu)
        \frac{g^2T(r)}{48\pi^2\epsilon}.
  \]
  It is canceled by
  \[
    \boxed{\delta Z_3\big|_\Phi
      =-\frac{g^2T(r)}{48\pi^2\epsilon}
      =-\frac{g^2}{16\pi^2\epsilon}\frac{T(r)}3.}
  \]
  Since this matter loop changes neither \(Z_1\) nor \(Z_2\), the defining
  relation in the hint gives
  \[
    g=g_0\left[1+\frac12\delta Z_3\big|_\Phi+\cdots\right].
  \]
  Equivalently, in
  \[
    \beta(g)=-\frac{g^3}{16\pi^2}b_0+O(g^5),
  \]
  one complex scalar contributes
  \[
    \boxed{\Delta b_0=-\frac13T(r),\qquad
      \Delta\beta(g)=+\frac{g^3T(r)}{48\pi^2}.}
  \]
  The positive beta-function contribution is the expected screening by
  charged matter.  A real scalar, when the representation permits one,
  has half as many independent fields and therefore contributes half this
  amount.

  \SupplementarySolution{10}{The One-Loop Three-Gluon Counterterm}
  Put \(p+q+r=0\), and abbreviate the tree tensor by
  \[
    {\cal V}_{\mu\nu\rho}^{abc}(p,q,r)
      =g f^{abc}V_{\mu\nu\rho}(p,q,r),\qquad
    h\equiv\frac{g^2}{16\pi^2\epsilon}.
  \]
  It is useful first to make the graph calculation explicit.  Apart from
  the common continuation factor, the two fermion orientations give
  \[
  \begin{split}
    \Gamma_{\!F}^{abc}={}&-n_Fg^3\mu^{2\epsilon}
      \int\frac{d^d\ell}{(2\pi)^d}\,
      \frac{\operatorname{tr}_{D}\!\left[
        \gamma_\mu\slashed\ell\gamma_\nu(\slashed\ell+\slashed q)
        \gamma_\rho(\slashed\ell-\slashed p)\right]}
      {\ell^2(\ell+q)^2(\ell-p)^2}\,
      \operatorname{tr}_r(T^aT^bT^c)\\
    &+\bigl[(b,\nu,q)\leftrightarrow(c,\rho,r)\bigr].
  \end{split}
  \]
  The leading cubic power of \(\ell\) is odd and integrates to zero; the
  ultraviolet pole comes from the terms linear in the external momenta.
  The difference of the two color orderings is
  \[
    \operatorname{tr}_r(T^aT^bT^c)
      -\operatorname{tr}_r(T^aT^cT^b)
      =iT(r)f^{bca}.
  \]
  This both removes the symmetric trace and produces the tree color
  tensor.  The initial minus sign is the closed-fermion-loop sign.

  For the two ghost orientations one may route the three denominators as
  \(k_1=\ell\), \(k_2=\ell+q\), and \(k_3=\ell-p\).  A representative
  orientation is
  \[
    \Gamma_{\rm gh}^{abc}
      =-g^3\mu^{2\epsilon}
       f^{ade}f^{bef}f^{cfd}
       \int\frac{d^d\ell}{(2\pi)^d}\,
       \frac{(k_1)_\mu(k_2)_\nu(k_3)_\rho}
            {k_1^2k_2^2k_3^2},
  \]
  and the second is obtained by reversing the ghost flow.  Here too the
  overall minus sign is the closed-Grassmann-loop sign.

  To display the remaining four graphs compactly, let
  \(W_{\mu\nu\rho\sigma}^{abcd}\) denote the complete color--Lorentz tensor
  inside the square brackets of the four-gluon rule in the question.  The
  gluon triangle is
  \[
  \begin{split}
    \Gamma_{\triangle}^{abc}
     ={}&g^3\mu^{2\epsilon}
       f^{ade}f^{bef}f^{cfd}
       \int\frac{d^d\ell}{(2\pi)^d}\,
       \frac{{\cal N}_{\mu\nu\rho}}{k_1^2k_2^2k_3^2},\\
    {\cal N}_{\mu\nu\rho}={}&
       V_{\mu\alpha\beta}(p,-k_2,k_1)
       V_{\nu\gamma\alpha'}(q,-k_3,k_2)\\
      &{}\times V_{\rho\beta'\gamma'}(r,-k_1,k_3)
       \eta^{\alpha\alpha'}\eta^{\beta\beta'}
       \eta^{\gamma\gamma'}.
  \end{split}
  \]
  A representative swordfish, with external legs \(a,b\) on the
  four-gluon vertex and \(c\) on the three-gluon vertex, is
  \[
    \Gamma_{{\rm sf},(ab)|c}
      =\frac12g^3\mu^{2\epsilon}
       \int\frac{d^d\ell}{(2\pi)^d}\,
       \frac{
         W_{\mu\nu\alpha\beta}^{abde}
         f^{dec}V_{\alpha'\beta'\rho}
            (-\ell,\ell+r,-r)
         \eta^{\alpha\alpha'}\eta^{\beta\beta'}}
       {\ell^2(\ell+r)^2}.
  \]
  The other two graphs are the cyclic relabelings
  \((ab)|c\to(bc)|a\to(ca)|b\).  The factor \(1/2\) in each expression is
  essential: exchanging the two indistinguishable internal gluon lines is
  an automorphism of the graph.

  Feynman parametrization reduces all four expressions to the pole
  integrals
  \[
  \begin{aligned}
    \mu^{2\epsilon}\int\frac{d^dL}{(2\pi)^d}
       \frac1{(L^2+\Delta-i0)^2}
       &=\frac{i}{16\pi^2\epsilon}+O(1),\\
    \mu^{2\epsilon}\int\frac{d^dL}{(2\pi)^d}
       \frac{L_\alpha L_\beta}{(L^2+\Delta-i0)^3}
       &=\frac{i\eta_{\alpha\beta}}{64\pi^2\epsilon}+O(1).
  \end{aligned}
  \]
  In the second formula we used
  \(L_\alpha L_\beta\mapsto\eta_{\alpha\beta}L^2/d\).
  The color reductions needed in addition to the fermion trace are
  \[
    f^{ade}f^{bde}=C_A\delta^{ab},\qquad
    f^{ade}f^{bef}f^{cfd}=\frac12C_Af^{abc},
  \]
  together with the Jacobi identity.

  The Lorentz reduction is best audited in Bose-symmetric graph families:
  the two orientations in each directed loop and the three cyclic
  swordfish graphs must be summed before comparing with the full tree
  tensor.  Define \(k_{\cal G}\) by
  \[
    \Gamma_{\cal G}^{\rm pole}
      =-h\,k_{\cal G}\,{\cal V}_{\mu\nu\rho}^{abc},
    \qquad
    \delta Z_{1,3g}^{({\cal G})}=h\,k_{\cal G}.
  \]
  Direct contraction of the numerators above gives
  \[
  \begin{array}{c|c}
    \text{Bose-symmetric graph family}&k_{\cal G}\\ \hline
    \text{two fermion orientations}
      &-\dfrac43\,n_FT(r)\\[3pt]
    \text{two ghost orientations}
      &\dfrac1{24}C_A\\[3pt]
    \text{gluon triangle plus all three half-weight swordfish graphs}
      &\dfrac58C_A
  \end{array}
  \]
  For example, the fermion numerator uses
  \(\operatorname{tr}_{D}(\gamma_\mu\gamma_\alpha
  \gamma_\nu\gamma_\beta\gamma_\rho\gamma_\gamma)\);
  after symmetric integration its three metric contractions reconstruct
  the three terms of \(V_{\mu\nu\rho}\).  The ghost numerator supplies the
  \(C_A/24\) row.  In the pure-gauge family, terms not separately
  proportional to the Bose-symmetric tree tensor cancel between the
  triangle and the three cyclic swordfish graphs; retaining their
  \(1/2\) factors leaves \(5C_A/8\).  This is also a strong routing check:
  a shift of \(\ell\) changes no entry in dimensional regularization.

  Summing the rows gives
  \[
    \boxed{\delta Z_{1,3g}
      =\frac{g^2}{16\pi^2\epsilon}
       \left(\frac23C_A-\frac43n_FT(r)\right).}
  \]
  Finally, the counterterms supplied in the question give independently
  \[
  \begin{split}
    \delta Z_1+\delta Z_3-\delta Z_2
      &=
      h\left[-C_2(r)-C_A+\frac53C_A
        -\frac43n_FT(r)+C_2(r)\right]\\
      &=h\left(\frac23C_A-\frac43n_FT(r)\right)
       =\delta Z_{1,3g}.
  \end{split}
  \]
  Thus the \(C_2(r)\) terms cancel and the explicit diagrams verify
  \[
    \boxed{Z_{1,3g}=\frac{Z_3Z_1}{Z_2}}
  \]
  at one loop.  Omitting either a reversed directed loop, any external-leg
  permutation, or the swordfish symmetry factor destroys this equality.

  \SupplementarySolution{11}{Schwinger--Dyson, Ward--Takahashi, and QED Renormalization}
  \begin{enumerate}
    \item[(a)]
    Under a constant phase,
    \(\psi\mapsto e^{i\alpha}\psi\) and
    \(\overline\psi\mapsto\overline\psi e^{-i\alpha}\).
    The covariant kinetic term and mass are invariant, and the gauge field
    is unchanged.  Promoting \(\alpha\) to \(\alpha(x)\) changes the
    fermion action by
    \[
      \delta I=-\int d^4x\,(\partial_\mu\alpha)j^\mu
        =\int d^4x\,\alpha\,\partial_\mu j^\mu
    \]
    up to a boundary term.  The Euler--Lagrange equations therefore give
    \(\partial_\mu j^\mu=0\).
    \item[(b)]
    Use the same local phase as a change of integration variables in
    \(\langle T\psi(x_1)\overline\psi(x_2)\rangle\).  The vector measure has
    unit Jacobian.  The action variation inserts
    \(\partial_\mu j^\mu(x)\), while the two external fields vary with
    opposite signs.  Since the integral itself cannot change, the
    coefficient of arbitrary \(\alpha(x)\) gives
    \[
      \boxed{\partial_\mu
      \langle Tj^\mu(x)\psi(x_1)\overline\psi(x_2)\rangle
       =-\bigl[\delta^{(4)}(x-x_1)-\delta^{(4)}(x-x_2)\bigr]
        \langle T\psi(x_1)\overline\psi(x_2)\rangle.}
    \]
    This is the Euclidean-sign convention of the source question.
    In a Minkowski convention in which the current insertion is obtained
    by differentiating \(e^{iI}\), both sides acquire the same overall
    factor of \(i\).
    \item[(c)]
    Translation invariance turns the three-point function into
    \(S(p+q)\Gamma^\mu(p+q,p)S(p)\).  Fourier transforming the derivative
    supplies \(q_\mu\); the two contact terms remove respectively the
    propagator on the left and right.  Multiplying by the inverse
    propagators therefore yields
    \[
      \boxed{q_\mu\Gamma^\mu(p+q,p)
        =S^{-1}(p+q)-S^{-1}(p).}
    \]
    Dividing by \(q_\mu\) in the regular zero-momentum limit gives
    \[
      \boxed{\Gamma^\mu(p,p)
        =\frac{\partial S^{-1}(p)}{\partial p_\mu}.}
    \]
    \item[(d)]
    Write \(S_R=Z_2^{-1}S_0\) and
    \(\Gamma_R^\mu=Z_1\Gamma_0^\mu\).  The bare and renormalized
    Ward--Takahashi identities have the same normalization only if
    \[
      \boxed{Z_1=Z_2}.
    \]
    Since
    \(e_0=\mu^\epsilon Z_1Z_2^{-1}Z_3^{-1/2}e\), it follows that charge
    renormalization is controlled by the photon factor \(Z_3\).

    No loop expansion was used in deriving the functional identity.
    Subject to a regulator preserving vector gauge symmetry, current
    conservation, the Ward--Takahashi identity, and \(Z_1=Z_2\) are
    nonperturbative.  Any truncation of Schwinger--Dyson equations must
    preserve this identity to avoid gauge-dependent charges.
  \end{enumerate}

\end{enumerate}
